% GENERATED FILE. Edit canonical lessons, then run npm run generate:books.
% canonical-source-hash: fnv1a64:18ff62351dd4d1bc
% canonical-lessons: BN-C287-muhurto, BN-C287-chokite, BN-C287-akre-dhora, BN-C287-totkkhonat, BN-C287-achhre-pora

\chapter{In a Split Second}
\label{ch:bn-in-a-split-second}
\hlchaptermodality{287}

\begin{chapteropening}
\textbf{By the end of this chapter:} \emph{I can tell what happened in one sudden moment.}

Complete the last lesson of chapter 287: I can tell what happened in one sudden moment.
\end{chapteropening}

\section[\texorpdfstring{\bn{মুহূর্ত} (muhūrtô)}{মুহূর্ত (muhūrtô)}]{\bn{মুহূর্ত} (muhūrtô) — a moment; an instant}

\label{lesson:BN-C287-muhurto}



\begin{quote}
Before the new one: say the Bengali for a shadow, then the Bengali for to observe.
\end{quote}

\subsection*{You'll want to know: \bn{মুহূর্ত}}
\textbf{\bn{মুহূর্ত}} — \emph{muhūrtô} — \textquotedblleft{}a moment; an instant\textquotedblright{}.

\begin{quote}
\bn{এক} \bn{মুহূর্তের} \bn{জন্য} \bn{আমি} \bn{কিছু} \bn{বুঝলাম} \bn{না।}
\end{quote}

\begin{quote}
\emph{ek muhūrter jônyô āmi kichhu bujhlām nā.}
\end{quote}

\begin{quote}
For a moment I didn't understand anything.
\end{quote}

\begin{cousinweb}
\textbf{\bn{মুহূর্ত}} (\emph{muhūrtô}) is Sanskrit \emph{muhūrta}, 'a moment, an instant' (Monier-Williams), and also a fixed thirtieth of the day; that is why the word also names the auspicious hour for a wedding.
\end{cousinweb}

\begin{grammarlens}[title={What you've built}]
One more word for telling what happened.
\end{grammarlens}

\subsection*{Your turn}
\begin{itemize}
\raggedright
  \item \emph{Say it:} \emph{muhūrtô}
  \item \emph{Say it:} \emph{muhūrtô}, once more
  \item \emph{Say it:} \emph{ek muhūrter jônyô āmi kichhu bujhlām nā.} — For a moment I didn't understand anything.
  \item \emph{Recall:} say the Bengali for a shadow, then the Bengali for to observe, then say \emph{muhūrtô} again
\end{itemize}

\begin{tcolorbox}[breakable,colback=teal!4,colframe=teal!35!black,title={Before you move on}]
What does \bn{মুহূর্ত} mean? (\textquotedblleft{}A moment; an instant\textquotedblright{}.) Say it once more.
\end{tcolorbox}
\section[\texorpdfstring{\bn{চকিতে} (chôkite)}{চকিতে (chôkite)}]{\bn{চকিতে} (chôkite) — in a flash; in an instant}

\label{lesson:BN-C287-chokite}



\begin{quote}
Before the new one: say the Bengali for to observe, then the Bengali for a moment.
\end{quote}

\subsection*{You'll want to know: \bn{চকিতে}}
\textbf{\bn{চকিতে}} — \emph{chôkite} — \textquotedblleft{}in a flash; in an instant\textquotedblright{}.

\begin{quote}
\bn{চকিতে} \bn{চোরটা} \bn{জানালা} \bn{দিয়ে} \bn{লাফিয়ে} \bn{পালাল।}
\end{quote}

\begin{quote}
\emph{chôkite chorṭā jānālā diye lāphiye pālālô.}
\end{quote}

\begin{quote}
In a flash the thief jumped out of the window and ran.
\end{quote}

\begin{cousinweb}
\textbf{\bn{চকিতে}} (\emph{chôkite}) is Sanskrit \emph{cakita}, 'trembling, startled' (Monier-Williams), in the -e form: 'in a start'. The sense 'in a flash' is the Bengali one.
\end{cousinweb}

\begin{grammarlens}[title={What you've built}]
One more word for telling what happened.
\end{grammarlens}

\subsection*{Your turn}
\begin{itemize}
\raggedright
  \item \emph{Say it:} \emph{chôkite}
  \item \emph{Say it:} \emph{chôkite}, once more
  \item \emph{Say it:} \emph{chôkite chorṭā jānālā diye lāphiye pālālô.} — In a flash the thief jumped out of the window and ran.
  \item \emph{Recall:} say the Bengali for to observe, then the Bengali for a moment, then say \emph{chôkite} again
\end{itemize}

\begin{tcolorbox}[breakable,colback=teal!4,colframe=teal!35!black,title={Before you move on}]
What does \bn{চকিতে} mean? (\textquotedblleft{}In a flash; in an instant\textquotedblright{}.) Say it once more.
\end{tcolorbox}
\section[\texorpdfstring{\bn{আঁকড়ে} \bn{ধরা} (ā̃kṛe dhôrā)}{আঁকড়ে ধরা (ā̃kṛe dhôrā)}]{\bn{আঁকড়ে} \bn{ধরা} (ā̃kṛe dhôrā) — to clutch; to cling to}

\label{lesson:BN-C287-akre-dhora}



\begin{quote}
Before the new one: say the Bengali for a moment, then the Bengali for in an instant.
\end{quote}

\subsection*{You'll want to know: \bn{আঁকড়ে} \bn{ধরা}}
\textbf{\bn{আঁকড়ে} \bn{ধরা}} — \emph{ā̃kṛe dhôrā} — \textquotedblleft{}to clutch; to cling to\textquotedblright{}.

\begin{quote}
\bn{ভয়ে} \bn{শিশুটি} \bn{বাবার} \bn{হাত} \bn{আঁকড়ে} \bn{ধরল।}
\end{quote}

\begin{quote}
\emph{bhôye shishuṭi bābār hāt ā̃kṛe dhôrlô.}
\end{quote}

\begin{quote}
The frightened child clung to his father's hand.
\end{quote}

Past (he/she …): \bn{সেই} \bn{শিশু} \bn{দড়িটা} \bn{আঁকড়ে} \bn{ধরল}, \emph{sei shishu dôṛiṭā ā̃kṛe dhôrlô}.

\begin{cousinweb}
\textbf{\bn{আঁকড়ে} \bn{ধরা}} (\emph{ā̃kṛe dhôrā}) is the -e form of \bn{আঁকড়া} — CDIAL 100 lists Bengali \emph{ã̄kṛā}, 'a hook, a tendril', under Sanskrit \emph{aṅka}, 'a hook' — before \textbf{\bn{ধরা}}, to hold: to hold on the way a hook holds.
\end{cousinweb}

\begin{grammarlens}[title={What you've built}]
One more word for telling what happened.
\end{grammarlens}

\subsection*{Your turn}
\begin{itemize}
\raggedright
  \item \emph{Say it:} \emph{ā̃kṛe dhôrā}
  \item \emph{Say it:} \emph{ā̃kṛe dhôrā}, once more
  \item \emph{Say it:} \emph{bhôye shishuṭi bābār hāt ā̃kṛe dhôrlô.} — The frightened child clung to his father's hand.
  \item \emph{Recall:} say the Bengali for a moment, then the Bengali for in an instant, then say \emph{ā̃kṛe dhôrā} again
\end{itemize}

\begin{tcolorbox}[breakable,colback=teal!4,colframe=teal!35!black,title={Before you move on}]
What does \bn{আঁকড়ে} \bn{ধরা} mean? (\textquotedblleft{}To clutch; to cling to\textquotedblright{}.) Say it once more.
\end{tcolorbox}
\section[\texorpdfstring{\bn{তৎক্ষণাৎ} (tôtkkhônāt)}{তৎক্ষণাৎ (tôtkkhônāt)}]{\bn{তৎক্ষণাৎ} (tôtkkhônāt) — at once; immediately}

\label{lesson:BN-C287-totkkhonat}



\begin{quote}
Before the new one: say the Bengali for in an instant, then the Bengali for to clutch.
\end{quote}

\subsection*{You'll want to know: \bn{তৎক্ষণাৎ}}
\textbf{\bn{তৎক্ষণাৎ}} — \emph{tôtkkhônāt} — \textquotedblleft{}at once; immediately\textquotedblright{}.

\begin{quote}
\bn{আওয়াজ} \bn{শুনে} \bn{দারোগা} \bn{তৎক্ষণাৎ} \bn{বাইরে} \bn{এলেন।}
\end{quote}

\begin{quote}
\emph{āoyāj shune dārogā tôtkkhônāt bāire elen.}
\end{quote}

\begin{quote}
Hearing the noise, the inspector came out at once.
\end{quote}

\begin{cousinweb}
\textbf{\bn{তৎক্ষণাৎ}} (\emph{tôtkkhônāt}) is Sanskrit \emph{tatkṣaṇāt}, 'at that very moment' (Monier-Williams): \emph{tat}, 'that', and \emph{kṣaṇa}, 'an instant'. It belongs to the news and to careful speech.
\end{cousinweb}

\begin{grammarlens}[title={What you've built}]
One more word for telling what happened.
\end{grammarlens}

\subsection*{Your turn}
\begin{itemize}
\raggedright
  \item \emph{Say it:} \emph{tôtkkhônāt}
  \item \emph{Say it:} \emph{tôtkkhônāt}, once more
  \item \emph{Say it:} \emph{āoyāj shune dārogā tôtkkhônāt bāire elen.} — Hearing the noise, the inspector came out at once.
  \item \emph{Recall:} say the Bengali for in an instant, then the Bengali for to clutch, then say \emph{tôtkkhônāt} again
\end{itemize}

\begin{tcolorbox}[breakable,colback=teal!4,colframe=teal!35!black,title={Before you move on}]
What does \bn{তৎক্ষণাৎ} mean? (\textquotedblleft{}At once; immediately\textquotedblright{}.) Say it once more.
\end{tcolorbox}
\section[\texorpdfstring{\bn{আছড়ে} \bn{পড়া} (āchhṛe pôṛā)}{আছড়ে পড়া (āchhṛe pôṛā)}]{\bn{আছড়ে} \bn{পড়া} (āchhṛe pôṛā) — to crash down; to be dashed down}

\label{lesson:BN-C287-achhre-pora}



\begin{quote}
Before the new one: say the Bengali for to clutch, then the Bengali for at once.
\end{quote}

\subsection*{You'll want to know: \bn{আছড়ে} \bn{পড়া}}
\textbf{\bn{আছড়ে} \bn{পড়া}} — \emph{āchhṛe pôṛā} — \textquotedblleft{}to crash down; to be dashed down\textquotedblright{}.

\begin{quote}
\bn{ঝড়ে} \bn{একটা} \bn{বড়} \bn{গাছ} \bn{বাড়ির} \bn{ছাদে} \bn{আছড়ে} \bn{পড়ল।}
\end{quote}

\begin{quote}
\emph{jhôṛe ekṭā bôṛo gāchh bāṛir chhāde āchhṛe pôṛlô.}
\end{quote}

\begin{quote}
In the storm a big tree came crashing down on the roof of the house.
\end{quote}

Past (he/she …): \bn{সেই} \bn{গাছ} \bn{রাস্তায়} \bn{আছড়ে} \bn{পড়ল}, \emph{sei gāchh rāstāy āchhṛe pôṛlô}.

\begin{cousinweb}
\textbf{\bn{আছড়ে} \bn{পড়া}} (\emph{āchhṛe pôṛā}) is the -e form of \bn{আছড়ানো}, to dash down — CDIAL 1081 lists Bengali \emph{āchṛāna}, 'to throw down hard', and \emph{āchāṛ}, 'a dashing down', under \emph{ācchardati}, 'pours on'; how the sense moved is not set out there — before \bn{পড়া} in its sense 'to fall'.
\end{cousinweb}

\begin{grammarlens}[title={What you've built}]
One more word for telling what happened.
\end{grammarlens}

\subsection*{Your turn}
\begin{itemize}
\raggedright
  \item \emph{Say it:} \emph{āchhṛe pôṛā}
  \item \emph{Say it:} \emph{āchhṛe pôṛā}, once more
  \item \emph{Say it:} \emph{jhôṛe ekṭā bôṛo gāchh bāṛir chhāde āchhṛe pôṛlô.} — In the storm a big tree came crashing down on the roof of the house.
  \item \emph{Recall:} say the Bengali for to clutch, then the Bengali for at once, then say \emph{āchhṛe pôṛā} again
\end{itemize}

\begin{tcolorbox}[breakable,colback=teal!4,colframe=teal!35!black,title={Before you move on}]
What does \bn{আছড়ে} \bn{পড়া} mean? (\textquotedblleft{}To crash down; to be dashed down\textquotedblright{}.) Say it once more.
\end{tcolorbox}
